\documentclass[a4paper]{article}

% Basic packages
\usepackage[fontset=fandol]{ctex}
\usepackage{amsmath, amssymb}
\usepackage{graphicx}
\usepackage[margin=2.5cm]{geometry}
\usepackage[most]{tcolorbox}
\usepackage{etoolbox}
\usepackage{listings}
\usepackage{hyperref}
\usepackage{booktabs}
\usepackage{subcaption}
\usepackage{float}
\usepackage{tikz}
\IfFileExists{pgfplots.sty}{
    \usepackage{pgfplots}
    \pgfplotsset{compat=1.18}
}{}

% Highlight boxes
\newtcolorbox{knowledgebox}[1]{
    enhanced, colback=blue!5!white, colframe=blue!75!black, colbacktitle=blue!75!black,
    coltitle=white, fonttitle=\bfseries, title=#1, attach boxed title to top left={yshift=-2mm, xshift=2mm},
    boxrule=1pt, sharp corners
}

\newtcolorbox{importantbox}[1]{
    enhanced, colback=yellow!10!white, colframe=yellow!80!black, colbacktitle=yellow!80!black,
    coltitle=black, fonttitle=\bfseries, title=#1, sharp corners
}

\newtcolorbox{warningbox}[1]{
    enhanced, colback=red!5!white, colframe=red!75!black, colbacktitle=red!75!black,
    coltitle=white, fonttitle=\bfseries, title=#1, sharp corners
}

\newtcolorbox{dialoguebox}[1]{
    enhanced, breakable, colback=green!4!white, colframe=green!45!black, colbacktitle=green!45!black,
    coltitle=white, fonttitle=\bfseries, title=#1, sharp corners
}

% Listings
\lstset{
    language=Python,
    basicstyle=\ttfamily\small,
    keywordstyle=\color{blue},
    stringstyle=\color{red!60!black},
    commentstyle=\color{green!60!black},
    breaklines=true,
    frame=single,
    numbers=left,
    numberstyle=\tiny\color{gray},
    captionpos=b,
    extendedchars=false
}

% Front-page metadata
\newcommand{\notetitle}{课程笔记：构建 OpenHands 的经验与教训}
\newcommand{\noteauthors}{Bowen-mantou（整理）}
\newcommand{\notedate}{2026 年 10 月 23 日}
\newcommand{\videochannel}{常青藤名校课（B 站）}
\newcommand{\videopublishdate}{2026-09-09}
\newcommand{\videoduration}{69:20}
\newcommand{\videourl}{https://www.bilibili.com/video/BV1a6Yx62EH4?p=28}
\newcommand{\videocoverpath}{cover.jpg}

\begin{document}

\begin{titlepage}
\centering
\vspace{1.2cm}
{\huge\bfseries \notetitle\par}
\vspace{0.8cm}
{\large \noteauthors\par}
\vspace{0.3cm}
{\large CMU 11-768: AI Agents（2026 年秋季）\par}
{\large 主讲：Graham Neubig\par}
\vspace{0.3cm}
{\large \notedate\par}
\vspace{1.2cm}

\ifdefempty{\videocoverpath}{
{\small 视频封面因版权未收录，见下方视频链接。\par}
}{
\par
}

\vfill
\begin{tcolorbox}[width=0.9\textwidth, colback=black!2!white, colframe=black!60, sharp corners]
\textbf{视频作者/频道}：\videochannel\par
\textbf{发布日期}：\videopublishdate\par
\textbf{视频时长}：\videoduration\par
\textbf{视频链接}：
\ifdefempty{\videourl}{
未填写
}{
\href{\videourl}{\nolinkurl{\videourl}}
}
\par\textbf{PDF制作Skill}：\href{https://github.com/wdkns/wdkns-skills}{\nolinkurl{https://github.com/wdkns/wdkns-skills}}
\end{tcolorbox}
\end{titlepage}

\tableofcontents
\newpage

%% --- 正文内容开始 --- %%

\section{从研究到产品：这一讲为什么不同}

\subsection{一堂关于“真正被用户使用”的 agent 课}

这是 Graham Neubig 在本课程中的最后一讲（后面是 Daniel Fried 的课与客座讲座）。开场他就说明，这一讲会和前面不太一样：一部分讲技术，但更多是讲他自己的经验——如何把一个 agent 产品真正做成有人用的东西，以及如何把“让用户真正用起来”的过程中遇到的痛点转化为研究课题。

这个定位很值得留意。前面十几讲大多在讲方法：怎么定义 agent、怎么设计 harness、怎么做评测、怎么做 RL 训练。而这一讲回答的是另一个问题：\textbf{当你把一个 agent 系统交给真实用户、放进真实工作流之后，哪些工程细节和研究问题才会暴露出来？}

\subsection{OpenHands 是什么：两条工作流}

OpenHands 是较早的开源软件开发 agent 之一，从 OpenDevin 改名而来。课上先演示了它的前端，核心是两条不同的工作流：

\begin{itemize}
\item \textbf{Human-in-the-loop}：在聊天窗口里输入你今天想做的任务，agent 在一个可交互的工作区里完成它；
\item \textbf{Automation（自动化）}：把重复性工作配置成自动化任务，比如代码评审（code review）、仓库监控、Slack 频道监控等，让 agent 在后台持续执行。
\end{itemize}

演示里，Neubig 让 agent 去自己的开源仓库里找“今天需要处理的 open PR”，agent 会创建一段对话、访问仓库、读取 PR 信息并给出待办列表。他还展示了自动化面板：一个“每日工作流准备”的 automation 会把今天要做的事读出来；一个“软件工厂”式的 automation 集合里包含 PR 评审、issue 分诊（triage）、Slack 监控、依赖失败自动修复等。这两条工作流对应两种完全不同的产品形态，后面会分别展开。

\subsection{与 Claude Code、Codex 的关键差异}

Neubig 点出了 OpenHands 与大家熟悉的 Claude Code、Codex 的核心区别：后两者主要绑定在自己的语言模型生态里（分别偏向 Claude 和 GPT 模型），而 OpenHands 从设计上就尽量做到\textbf{模型无关（model-agnostic）}：它支持很多不同模型，可以临时切换模型再切回来，也可以自定义 MCP server、skills、plugins、apps。

另一个差异在“怎么用它”。OpenHands 是 Python 写的、面向“直接写代码调用它”的场景；OpenCode、Pi 等开源 agent 多用 TypeScript，且很多工具主要围绕 CLI 交互（很多 benchmark 就是通过命令行参数驱动 agent 跑起来的）。用 CLI 调用很方便，跨语言也能用，但一旦你想深入改造、实现新模块，命令行参数和配置文件的表达能力就会成为限制。因此选型时可以先问自己：你更想“当命令行工具用”，还是“当 Python 库改”？

\begin{knowledgebox}{Pi 的哲学：裸核 + 扩展}
Neubig 顺带对比了 Pi 的设计思路：Pi 刻意只给一个非常精简的内核，然后靠扩展把功能加上去。OpenHands 的 SDK 走了相反的路：先给一套开箱即用的好默认值，但允许你把任何一部分拆掉替换。两种路线没有绝对优劣，反映的是“最小内核”与“默认体验”的优先级差异。
\end{knowledgebox}

\subsection{本章小结}

\begin{itemize}
\item 这一讲的主题是“把 agent 做成产品”的经验：真实用户、真实工作流会暴露研究论文里看不到的问题。
\item OpenHands 支持两条工作流：交互式的 human-in-the-loop，以及后台自动化的 automation；它的定位是模型无关、可用 Python 直接编程的 agent 平台。
\item 与 Claude Code、Codex、OpenCode、Pi 等工具的差异，不在“谁更聪明”，而在模型生态绑定、语言与接口形态、可改造程度上。
\end{itemize}

\section{历史：从代码补全到 agentic coding}

\subsection{十年代码生成，两年 agent}

Neubig 从 2016 年就开始做代码生成，但直到 2023 年前后，主流做法仍然不是 agentic 的：主要是代码补全、生成单个函数。真正的 agentic code generation 大约从 2024 年才开始，而短短两年里，他自己已经不再手写代码了。

2023 年的一件大事是 SWE-bench 的发布：它把真实 GitHub issue、仓库快照、patch 和仓库测试组合在一起，成为一个仓库级的评测基准。它的影响力极大，但早期结果很差——大约 5\% 的水平，几乎没人能拿到像样的分数。这个“看起来能做、实际做不出”的落差，正是 agent 研究要解决的问题。

\subsection{Devin 时刻与开源复现}

真正让局面改变的，是 2024 年 3 月 Devin 的演示。Devin 被称为“第一个 AI 软件工程师”：你输入任务，它会自己规划、开始工作。这段演示视频在很短时间内获得上百万播放。Neubig 当时已经做了几年 web agent 与 coding agent，也有八年代码生成的积累；他看完的第一反应是“我知道这件事是可能的”，但 Devin 把它包装得非常有吸引力。

问题在于 Devin 不开源。Neubig 一直是开源的支持者，于是萌生了“做一个类似但开源的东西”的想法。就在 Devin 演示的第二天，Qwen 团队的 Junyang Lin 与 Binyuan Hui 发出号召：一起复现 Devin。于是这个只有 README 的仓库迅速拿到一万 star，成为当时最受关注的项目之一。

Neubig 当天晚上用 Claude 的聊天窗口，照着一张 Devin 的截图，做出了第一版界面；大约三天后，就出现了第一个真正能跑起来、能完成任务的 agent 例子（OpenDevin，后来改名 OpenHands）。与此同时，SWE-bench 的作者们其实早就在做 coding agent，后来发布了 SWE-agent；只是发布节奏被 Devin 抢了一点先机。5 月，OpenDevin 发布 CodeAct 1.0，把“写代码执行”作为核心动作空间。整个领域在几个月内快速成形。

\subsection{什么时候 agent 才真正可用}

但必须诚实地说：在很长一段时间里，这些 agent 并不真的可用。Neubig 提到，Devin 团队自己也承认这一点——他们甚至在公交站打广告说“Devin, it's actually good now”，这本身就说明之前不够好。从 2024 年到 2024 年 11 月左右，agent 的整体表现都很勉强。

转折点在 2024 年年底：Claude Sonnet 3.5/3.7 这一代模型让 agent 开始“真的能用”；2025 年则是 coding agent 全面变好的一年，到年底，不只是 Anthropic，OpenAI 等大公司也都有了很强的 coding 模型。到了 2026 年，Neubig 的判断是：模型还在变强，但这一年的主战场是\textbf{部署与界面}——更长的运行时间、并行任务、更好的可视化、可重复的工作流、凭证管理与验证。

这段历史的用处不在于怀旧，而在于解释后面要讲的设计：很多看似“产品细节”的东西（比如上下文压缩、skills、automation、可验证性），都是过去一年半里从真实使用中被逼出来的经验，后来才被固化进框架。

\subsection{本章小结}

\begin{itemize}
\item 代码生成从 2016 年就开始了，但 agentic 形态是 2023—2024 年才出现的；SWE-bench（2023.10）用真实 issue 与仓库测试定义了“仓库级”问题，早期成绩只有约 5\%。
\item 2024 年 3 月 Devin 演示引爆关注，次日 Qwen 团队号召复现，OpenDevin（后为 OpenHands）在几天内跑通；SWE-agent（4 月）与 CodeAct 1.0（5 月）奠定了早期动作空间。
\item 2024 年底之前 agent 普遍不好用；2024 年底到 2025 年模型变强后 agent 才真正可用；2026 年的关键在部署与界面。
\end{itemize}

\section{软件 agent SDK：harness 的设计与教训}

\subsection{第二代 harness：好默认值，但全部可拆}

OpenHands 的 harness 已经发展到第二代。Neubig 很坦白：第一代 harness 很糟糕，难用、难扩展。大约八个月前，他们干脆把原有实现全部推倒重来，围绕“可组合、可替换”重新设计。

这个 SDK 的设计哲学是：给一套真的很好用的默认值，开箱即用；但允许你把它们拆掉，替换成自己的实现。它覆盖 models、tools、state、workspace 等基础件，也允许自定义 model choice、新工具、instructions 和部署方式。GUI、CLI 和各种集成（integrations）都建立在这个 SDK 之上。

下面这段代码就是课上展示的最小配置思路（按幻灯片内容整理，省略了无关细节）：

\begin{lstlisting}[caption={用 OpenHands SDK 配置一个基础 agent（按幻灯片内容整理）},label={lst:sdk-config}]
llm = LLM(
    model="openai/gpt-5.5",
    api_key=os.getenv("LLM_API_KEY"),
)
agent = Agent(
    llm=llm,
    tools=[
        Tool(name=TerminalTool.name),
        Tool(name=FileEditorTool.name),
        Tool(name=TaskTrackerTool.name),
    ],
)
\end{lstlisting}

配置好 agent 之后，workflow 也很直接：指定 workspace（工作目录），创建一段 conversation，发消息，然后运行。这里 Neubig 特别提醒：\textbf{workspace 是灵活的}——你可以在本地目录、容器、远端沙箱里跑，这取决于你怎么做隔离。而 conversation 的本质是 tasks、actions、observations 三者的累积，再加上一个 action-observation 循环。

\subsection{自定义工具：action、observation 与 executor}

SDK 允许你实现自定义工具。做法是把“动作”和“观测”都定义成类型：例如实现一个 grep 工具，就定义 \texttt{GrepAction} 与 \texttt{GrepObservation}，再提供一个 executor——它接收一个 action，执行后返回一个 observation。这相当于给“工具”定义了一个通用接口：动作与观测是数据，executor 是把数据变成环境行为的函数。

\begin{lstlisting}[caption={自定义工具：action 与 observation 的类型定义（按幻灯片内容整理）},label={lst:custom-tool}]
class GrepAction(Action):
    pattern: str = Field(description="Pattern to search for")
    path: str = Field(default=".")

class GrepObservation(Observation):
    matches: list[str] = Field(default_factory=list)
\end{lstlisting}

这段设计看起来简单，但它把“工具”从一个黑盒函数变成了可记录、可回放、可比较的数据结构——这正是后面做评测、做训练、做自动化时需要的东西。

\subsection{隔离与安全：沙箱、确认策略与安全分析器}

Agent 会执行任意命令，所以隔离是绕不开的。课上讲了两种沙箱用法：一种是先准备好 sandbox，再把 agent 跑在里面；另一种是使用 Python 的 \texttt{with} 语句，让 with 块内执行的所有命令都进入沙箱。Neubig 说两种都不错，他没有明确结论，但“直接把 agent 跑在沙箱里”的生态工具更多一些。

除了沙箱，还有一层安全机制是\textbf{确认策略（confirmation policy）}：决定哪些动作可以自主执行，哪些必须停下来等人确认；以及\textbf{安全分析器（security analyzer）}：在每一步分析这个动作是否安全到可以自主执行。前者是流程控制，后者是风险判断，两者结合才能在“自主性”和“可控性”之间取平衡。

SDK 的完整能力还包括 secrets 管理、LLM metrics、LLM routing、planning agent、observability 等。Neubig 的提醒是：这个 SDK 非常全，全到课上只能讲很小一部分；如果你需要长 rollout、内建可观测性等能力，它很合适。

\subsection{上下文管理：condensation（后来被改叫 compaction）}

长会话必然要压缩上下文。Neubig 说他们很早就实现了这件事，最初叫 \textbf{condensation}（压缩/凝缩），后来 Anthropic 把同一件事叫作 \textbf{compaction}，于是业界大多改用了 compaction 这个名字，而 OpenHands 内部仍叫 condensation。

关于压缩，有几个值得记住的工程经验：

\begin{itemize}
\item 最直接的实现是“旧事件 → 摘要”，并把\textbf{最初的指令}保留在近处（\texttt{keep\_first}），避免关键约束被摘要掉；
\item 他们试过把单个动作掩掉（mask out individual actions）的做法，但效果不如摘要，现在业界普遍采用摘要式压缩；
\item 何时压缩？可以每 $n$ 轮压缩一次，也可以按 token 数压缩；一个常用阈值是把模型最大上下文窗口乘以约 3/4，超过就触发压缩。
\end{itemize}

\subsection{研究选型建议}

针对“做研究该用哪个 harness”，Neubig 给了很实用的建议：如果你用 Python 做 coding agent 研究，可以考虑功能完整的 OpenHands，或者极简的 mini-swe-agent——后者只有 bash 工具、代码很少，方便你直接改。你要权衡的是：需要长 rollout、内建可观测性等“能力”，还是需要“能快速 hack 的简单底座”。如果你用 TypeScript，或者只通过 CLI 驱动，其他 agent 可能更顺手。

他还分享了两个常见问题的答案。为什么用 Python？因为项目最初是机器学习背景的人做的，而且 Python 在“读 agent 写的代码”时更快、更直观。代价是类型安全（需要用严格类型检查和每次提交跑类型检查来兜底）和速度，但速度通常不是瓶颈——你等的是模型或 agent 运行的程序，而不是 agent 框架本身。

\begin{knowledgebox}{模型会不会“只对 OpenHands 过拟合”？}
有同学问：模型是不是过拟合了 OpenHands 的结构？Neubig 的回答很坦然：确实有很多模型在 OpenHands 这类流行 harness 上训练过；他们也曾仔细研究 Claude 使用的工具集并尽量对齐（OpenHands 的工具与 Claude Code 很接近，因此和 Claude 模型配合很好）。但现在模型对新 harness 的适应能力已经不错，OpenAI 的模型对各种 harness 都能用；Gemini 曾经因为 harness 适配问题分数偏低，后来通过 multi-harness training 改善。总体结论是：主流模型现在大多开箱就能用。
\end{knowledgebox}

\subsection{本章小结}

\begin{itemize}
\item 第二代 OpenHands SDK 的设计原则是“好默认值 + 全部可拆换”；它把 models、tools、state、workspace、conversation 做成可组合的积木，十行左右代码就能搭出一个 agent。
\item 自定义工具 = action + observation + executor；conversation 是 tasks/actions/observations 的累积，workspace 决定隔离方式。
\item 工程化的安全与上下文管理是刚需：沙箱、确认策略、安全分析器、condensation（compaction）；压缩时要保留关键指令，阈值常用最大上下文窗口的约 3/4。
\end{itemize}

\section{单 agent 还是多 agent}

\subsection{2024 年的 specialized roles 风潮}

OpenHands 起步的时候，多智能体正流行。一个典型做法是给 agent 分配固定角色：Manager 负责计划、Reproducer 负责复现 bug、Fault localizer 负责定位、Editor 负责改代码、Verifier 负责验证修复。Neubig 提到这来自一篇多智能体工作（CodeR；他现场先口误说成 MetaGPT，随即更正），当时围绕 PR 的解决流程出现了大约 15 种变体。

\subsection{单 agent 哲学：Don't Sleep on Single Agent Systems}

面对这股风潮，Neubig 在 2024 年年底写了一篇博客《Don't Sleep on Single Agent Systems》，核心论证是：语言模型正在变强，而固定工作流里的角色分工，本质上是把“先做这个、再做那个”的流程写死；如果模型的指令遵循能力足够好，那么用一段提示词描述流程，模型自己就能完成，不一定需要复杂的多智能体编排。

所以从第一天起，他就拒绝在多智能体系统上投入太多精力。事后看，这个判断大概率是对的：至少对 coding 来说，现在已经很少有人实现那种固定的多智能体工作流了。他也留了例外：客服等场景可能需要严格固定的流程，因为一旦 agent 自由发挥就可能给出错误信息、带来法律责任；但在 coding 这种需要灵活性的任务上，单 agent 的收益通常更大。

\subsection{单 agent 的工程基础：通用工具、skills 与 MCP}

单 agent 能成立，前提是工具足够通用。OpenHands 给 agent 的主要能力只有三件：编辑文件、运行命令、浏览网页（视觉检查与文档查阅）。他们不额外堆很多专用工具，但支持 MCP（Model Context Protocol）把外部服务接进来。

但只有通用工具还不够：如果所有情形的操作规范都写进单一系统提示，提示词会膨胀到难以维护（“推 GitHub 时这么做、用 Docker 时那么做……”）。OpenHands 的解法是 \textbf{skills}：把特定场景的提示词（有时还包括脚本）打包成技能，在需要时拉进上下文。有趣的是，他们很早就有这个机制，当时叫 microagents，后来 Anthropic 实现了同样的东西并起名 skills——Neubig 承认，这个名字确实比 microagent 好得多。

两者还有一个小差别：Anthropic 的做法是让 agent 自己选择要查哪个 skill，再读取对应文件；OpenHands 则用显式的关键词触发。无论哪种，目标都是一致的：\textbf{在正确的时间，把正确的上下文拉进来}。

这个思路也给出了多智能体的替代方案：多智能体是给每个角色固定提示词和固定工具；skills + MCP 则是在需要时动态拉入提示词与工具。单个 agent 因此也能覆盖“角色分工”带来的能力。

\subsection{单 agent 能有多通用：OpenHands-Versa}

为了验证“通用工具 + 多模态浏览”的路线，OpenHands 团队（包括这门课的助教 Aditya）做了一篇论文《Coding Agents with Multimodal Browsing Are Generalist Problem Solvers》。思路很简单：给 agent 代码能力，再给它多模态浏览（看图像、看网页）的能力，它就能解很多不同类型的任务，而不只是写代码。

论文的核心结论是：这套通用组合不仅可行，而且在多个原本由专用 agent 主导的任务上超过了当时的最好结果。这成为 OpenHands 一以贯之的哲学主线：与其写一堆专用角色，不如把通用能力做强。

\subsection{多 agent 仍然值得做的三个理由}

不过，Neubig 并没有把多智能体一棍子打死。随着任务越来越有野心，多智能体在某些场景下越来越有用，他给出三个理由：

\begin{enumerate}
\item \textbf{效率}：如果任务彼此独立（例如十个不同的 PR），可以并行分发；
\item \textbf{独立评审}：让一个 agent 检查另一个 agent 的工作。这里有一个很有价值的经验观察：如果让 agent 评审自己的工作，它会宽容得多，因为它拥有一条条决策的“理由”；而一个看不到这些理由、上下文被切断的独立 agent，评审会严厉得多——这通常正是检查所需要的；
\item \textbf{成本控制}：把任务派给子 agent，它完成后只报告重要信息、丢掉过程中的上下文；也可以用更便宜的模型做实现，把更聪明的模型留给编排。
\end{enumerate}

\begin{warningbox}{多 agent 不是“白拿”的收益}
Neubig 特别强调，多智能体系统的上下文管理非常难。他们有一篇论文研究异步软件工程 agent：一个 manager 把任务分派给多个 engineer，在 PaperBench、Commit0 这类很长的实现任务上运行。要让它真正工作，需要工程师 agent 反馈“恰到好处”的上下文、避免彼此冲突或重复实现同一个东西，manager 还要把各人的产出正确集成回代码仓库——这些都很难。最后虽然做成了，但过程相当艰难。与其说多 agent 是捷径，不如说它把难点从“单 agent 的能力”转移到了“协调与集成”。
\end{warningbox}

课上还回答了同学的问题：多智能体怎么追踪 reasoning trace、怎么理解到底发生了什么？Neubig 的坦诚回答是：通常做得并不好；他自己的做法是直接问 manager，让 manager 去翻看 trace，因为它比自己快。多 agent 在 deep research 这类“把任务分发到大量文档源”的场景里更容易赢，但在需要精细集成的 coding 场景里，固定流程往往不值得。

\subsection{本章小结}

\begin{itemize}
\item 2024 年流行“固定角色 + 固定工作流”的多智能体；Neubig 主张单 agent：模型指令遵循变强后，一段提示词就能承载流程，coding 领域尤其如此。
\item 单 agent 的工程基础是通用工具（编辑/运行/浏览）+ skills（按需拉入场景提示词）+ MCP（按需接入外部工具），目标都是“在正确时间拉入正确上下文”。
\item OpenHands-Versa 论文说明通用能力可以跨任务取胜；多 agent 的价值主要在并行、独立评审、成本控制，但协调与集成成本很高，不是白拿的收益。
\end{itemize}

\section{Agent 跑在哪、怎么交互}

\subsection{本地还是云：一个并不显然的答案}

这是课上“实践部分”的第一个问题：agent 到底跑在你自己的电脑上，还是跑在云工作区里？

Devin 最初是云产品（SaaS）；最早的 OpenDevin 则跑在本地，但被隔离在 Docker 容器里——主要原因很现实：当时他们不信任 agent 直接在自己的机器上跑。Neubig 自己也曾认为本地模式不可能被接受：如果 agent 被 prompt injection，偷走你的数据或上传 API key 怎么办？

结果恰恰相反：人们非常喜欢在本地跑 agent。对人机协作型的日常工作（修 bug、实现功能），本地模式已经变成主流；而对非 coding 工作或后台工作（比如 code review），云工作区更流行。Claude Code 一类产品则在尝试混合形态：让 agent 跑在云上、却操作你本地的电脑。

\begin{warningbox}{本地执行的安全代价}
本地模式好用，但它把风险留给了用户：如果 agent 被 prompt injection 或权限过大，它可能泄露本地数据、凭证或修改不该改的文件。这也是为什么 OpenHands 强调沙箱、凭证限制与确认策略——安全不能靠“agent 应该乖”，而要靠运行时机制。
\end{warningbox}

\subsection{四种交互界面：IDE、GUI、CLI 与集成}

课上把当前主流的 agent 交互方式归为四类，并现场统计了同学们的使用情况：

\begin{table}[H]
\centering
\small
\caption{四类 agent 交互界面的特点（按课堂讲解整理）}
\begin{tabular}{p{0.14\textwidth}p{0.36\textwidth}p{0.40\textwidth}}
\toprule
界面 & 优势 & 局限或注意点 \\
\midrule
IDE & 熟悉：文件、diff、测试、调试器、聊天与任务管理、评审都在一个地方 & 如果你已经不怎么读代码，IDE 某种程度上是“遗留界面” \\
GUI & 为 agent 原生设计的桌面应用，界面可以围绕任务而非代码组织 & 使用习惯仍需培养，生态较新 \\
CLI & 到处都能用：服务器、家用机器、远端机器 & 交互形态偏文本，复杂可视化与评审体验有限 \\
集成 & 在已有工作流（Slack、GitHub、code review）里直接使用，不用切换界面 & 依赖平台能力，定制空间较小 \\
\bottomrule
\end{tabular}
\end{table}

课堂上 CLI 是使用最多的（约六成），IDE 次之（约一半略少），GUI 很少。Neubig 认为 CLI 的优势是“哪里都能用”，而 IDE 的问题在于：当 agent 承担了大部分编码工作、你不再逐行看代码时，IDE 的很多功能并非为你设计。GUI 的想象空间在于重新设计“以 agent 为主、代码为辅”的工作界面。集成式（如 Copilot code review）则可能越来越普遍——最好的界面是你已经待着的界面。

\subsection{人机协作的实证研究：更快，但更无聊、更难验证}

OpenHands 团队做过一项人机研究，比较普通 Copilot 用户与使用 OpenHands agent 的用户在 Python 任务上的表现。结果是“天壤之别”：agent 用户大约用一半时间完成了两倍的任务，正确率也显著更高（图中约 0.6 对 0.25，时间约 12 分钟对 25 分钟）。

但负面反馈同样重要：受试者说“看着 agent 干活很无聊”，以及“很难确认它到底做得好不好”。这两个吐槽恰好指出了人机交互设计的两个核心问题：

\begin{itemize}
\item 既然人会无聊，产品就应该支持多任务——让一个人同时推进多条 agent 工作流；
\item 既然验证困难，产品就必须帮助人验证输出——这不是“agent 能力”问题，而是界面与流程设计问题。
\end{itemize}

\subsection{本章小结}

\begin{itemize}
\item 本地 vs 云不是非此即彼：人机协作的日常 coding 偏本地，后台/非 coding 工作偏云；本地模式必须以沙箱与权限机制控制风险。
\item 四类界面各有取舍：IDE 熟悉但可能“遗留”，GUI 面向 agent 重新设计，CLI 无处不在，集成式把 agent 嵌入已有工作流；课堂使用率 CLI 最高、IDE 次之。
\item 人机研究显示 agent 显著提升效率（约一半时间、两倍任务、更高正确率），但“无聊”和“难验证”是两个必须由产品设计解决的问题。
\end{itemize}

\section{自动化与持续改进：把 agent 变成软件生产线}

\subsection{Automation 的形态：事件 → 有界任务 → 可检查的结果}

Neubig 说这是他们现在投入很多的方向：面向软件开发的自动化。一个 automation run 的结构可以概括为：一个触发事件（issue、PR 或定时器）进来，收集上下文，在隔离的工作区里运行，产出一个可检查的 artifact，最后由人或策略来检查。

具体到配置，一个 automation 需要：触发方式与指令、仓库与权限、输入事件到运行的映射，以及对输出质量的检查。他们在 OpenHands 里维护了一套“软件工厂”式的 automation 集合：GitHub PR 评审、issue 分诊、依赖版本升级失败自动修复、Slack 消息触发等。

\subsection{两个真实用例：issue 分诊与 PR 评审}

\paragraph{Issue triage。}作为一个大型开源项目，OpenHands 会收到大量 PR——其中很多是为了丰富简历而用 AI 快速生成的（“vibe coded”）。这未必是坏事：别人花 token 帮你实现功能；但麻烦在于很多 PR 做得不对，因为 issue 本身描述不清。于是 issue triage bot 做两件事：一是判断 issue 是否足够清楚，如果不够清楚，就列出需要澄清的问题；二是如果 issue 可做，就给出“完成条件/验收标准”（acceptance criteria），让人和 agent 都清楚要做到什么程度。

\paragraph{PR review。}他们有一套非常细的代码评审规则：比如这个 PR 是否在本仓库的范围内、需要检查哪些点、在什么情况下可以批准、什么情况下必须请求修改。这些规则进入 code review bot，由它在每个 PR 上执行。

为什么要花力气做这两件事？因为实现层面可以交给社区贡献者，而“判断该不该做”和“做得对不对”是维护者的瓶颈。更重要的是：这些流程每天重复，天然适合被自动化，也天然适合被改进。

\subsection{从历史运行中学习：Skill Refiner}

一旦 automation 每天运行，你就有了一个持续产生的数据集：每次 triage 或 review 都会产生一条轨迹，并且事后能拿到“延迟反馈”——比如 bot 说“LGTM（看起来没问题）”，但维护者随后说“这很差，需要修改”。于是问题就变成了一个机器学习问题：如何从这些带反馈的历史轨迹里，自动改进下一次的 prompt / skill？

这就是他们最近发表的 Skill Refiner 的工作。流程大致是：

\begin{enumerate}
\item 把轨迹按结果分成成功与失败，并附带证据（例如“维护者否定了 PR”就是失败证据）；
\item 用 embedding 对成功轨迹、失败轨迹分别聚类，找出反复出现的问题模式；
\item 从每个簇里生成对 skill（指令）的修改提案；
\item 用证据门控过滤：只保留那些有支持证据、看起来真能改善下一次运行的提案；
\item 把正、负提案合并，生成修订后的 skill。
\end{enumerate}

这比一些 baseline（课上提到 GEPA 与 T2S 等方法）更便宜也更好。

\subsection{下一步：从轨迹中挖掘工作流，以及严格评测}

Neubig 认为这个方向还有两块很大的空间：

\begin{itemize}
\item \textbf{从大量轨迹中挖掘 workflow}：现在的 automation 大多是人凭经验、声明式地写出来的；如果能聚类大量用户的使用轨迹，把常见用例自动变成 automation，再让它进入“自动改进”的闭环，系统就能持续变好；
\item \textbf{改 prompt 之外，还可以训练模型}：修改 prompt 或 harness 是一种改进，修改模型参数是另一种改进；后者需要更严格的评测。为 issue triage、code review 这类“欠定义（under-specified）”任务设计严格评测很难，但价值很大——如果你能把评测做好，就等于拿到了持续改进这些系统的钥匙。
\end{itemize}

\subsection{本章小结}

\begin{itemize}
\item Automation 把开发流程变成“事件 → 有界任务 → 可检查产物”；OpenHands 用它做 PR 评审、issue 分诊、依赖修复、Slack 触发等。
\item 重复运行的 automation 会积累带延迟反馈的轨迹，把“改进 agent”变成一个可以从数据中学习的问题；Skill Refiner 用聚类、提案与证据门控自动改进 skill，实验结果优于初始 skill 与 baseline。
\item 未来的两个方向：从轨迹中挖掘工作流，以及为欠定义任务建立严格评测（这也是训练模型的前提）。
\end{itemize}

\section{总结与延伸}

\subsection{讲者的收尾}

这是 Neubig 在本课程的最后一次课。他最后提到，之后会是客座讲座与 Daniel 的课程，并祝大家秋假愉快。整讲真正想传递的，是一个从开源项目到公司、从研究到产品的完整经验：

\begin{itemize}
\item 一个好的 agent 产品，往往不是从“最强模型”出发，而是从\textbf{真实工作流}出发：哪些步骤重复、哪些判断是瓶颈、哪些结果需要被检查；
\item 很多关键设计（上下文压缩、skills、sandbox、confirmation policy、automation）都是在真实使用中被迫形成的，后来才被抽象进框架；
\item 开源与公司化并不矛盾：OpenHands 变成公司，是因为它足够流行、有人愿意为它全职投入；而目前最大的客户是希望在自己的云上安全运行 coding agent、需要模型选择与可定制性的企业；
\item 研究选题可以从“用户痛点”中来：多智能体的协调成本、长任务的验证、欠定义任务的评测，都是真实系统里长出来的研究问题。
\end{itemize}

\subsection{全讲结构压缩}

把这一讲压缩成五句话：

\begin{enumerate}
\item \textbf{历史}：代码生成已有十年，但 agentic coding 只有两年；SWE-bench 定义了仓库级问题，Devin 引爆关注，开源复现与模型进步共同让 agent 在 2024 年底真正可用，2026 年的重点是部署与界面。
\item \textbf{harness}：好的 SDK 要给可用的默认值，同时允许全部拆换；工具 = action + observation + executor；会话 = tasks/actions/observations 的累积；沙箱、确认策略、安全分析器与上下文压缩（condensation/compaction）是工程刚需。
\item \textbf{单 agent 优先}：模型变强后，固定角色的多智能体工作流价值下降；通用工具 + skills + MCP 能在“正确时间拉入正确上下文”，单 agent 也能通用（OpenHands-Versa）。
\item \textbf{多 agent 有条件地用}：并行、独立评审、成本控制是三个正当理由；但协调与集成成本很高，独立评审甚至比自我评审更严厉，这是它真正的价值所在。
\item \textbf{自动化闭环}：automation 把开发流程变成可重复、可检查的生产线；历史轨迹 + 延迟反馈 = 可学习的改进信号（Skill Refiner）；最终目标是让 prompt、harness、模型与评测一起持续改进。
\end{enumerate}

\subsection{开放问题与实践建议}

\begin{itemize}
\item \textbf{长任务的验证}：agent 做得越久，人越难确认它是否真的完成；如何设计“可检查的产物”与验证流程，是产品与研究的共同问题。
\item \textbf{多智能体的协调}：如何让子 agent 的上下文“不多不少”、如何避免重复实现、如何把产出集成回仓库，仍是开放问题。
\item \textbf{欠定义任务的评测}：issue triage、code review 这类任务的“好”很难量化；谁能把评测做好，谁就能把自动化做得更好。
\item \textbf{从轨迹到工作流}：能否从真实使用轨迹中自动挖掘出 automation，而不是靠人凭经验手写？
\item \textbf{成本与安全}：多 agent 不一定更省钱，本地执行不一定更安全；两者都需要显式的优化目标与运行时约束。
\end{itemize}

如果要把这一讲变成可执行的一步，建议是：选一个你每天重复的软件工程流程（比如 PR 评审或 issue 分诊），把它写成一条 automation（触发 → 上下文 → 隔离执行 → 可检查产物），然后开始积累轨迹与反馈——这就是把 agent 从“demo”变成“系统”的最短路径。

\subsection{拓展阅读}

\begin{itemize}
\item OpenHands（开源软件开发 agent）：\url{https://github.com/All-Hands-AI/OpenHands}
\item OpenHands 文档：\url{https://docs.all-hands.dev/}
\item mini-swe-agent（极简 harness）：\url{https://github.com/SWE-agent/mini-swe-agent}
\item SWE-bench：\url{https://www.swebench.com/}
\end{itemize}

%% --- 正文内容结束 --- %%

\end{document}
